\chapter*{How to read this book}
\addcontentsline{toc}{chapter}{How to read this book}
\markboth{How to read this book}{How to read this book}

\section*{Who this book is for}

This book assumes you have completed roughly two years of an undergraduate
engineering or applied-mathematics curriculum. Concretely, we assume you are
comfortable with:

\begin{itemize}[nosep]
  \item \textbf{Vector calculus}---gradient, divergence, and curl; line, surface,
  and volume integrals; the divergence and Stokes theorems; integration by parts.
  \item \textbf{Linear algebra}---matrices and vectors, inner products and norms,
  orthogonality and projections, symmetric/positive-definite matrices,
  eigenvalues and eigenvectors, and the idea of solving $\mat{A}\vx=\vb$.
  \item \textbf{Ordinary differential equations}---linear systems, and the idea
  of marching a solution forward in time.
  \item \textbf{A first course in electromagnetism}---you have seen Maxwell's
  equations at least once, even if they are hazy.
\end{itemize}

We do \emph{not} assume any background in functional programming, type theory,
the finite-element method, or numerical iterative solvers. Those are precisely
what the book teaches. Where a topic needs more mathematics than the list above,
we build it up from the list above.

\section*{The shape of the book}

The book has six parts and a set of appendices, arranged so that each part rests
on the ones before it. \Cref{fig:roadmap} shows how they depend on one another;
you can read it as a suggested route.

\begin{figure}[t]
\centering
\begin{tikzpicture}[node distance=6mm and 10mm,
  pt/.style={stage, minimum width=34mm, minimum height=9mm},
  ap/.style={opbox, minimum width=34mm, minimum height=8mm, fill=simBgViolet, draw=simViolet}]
  \node[pt] (p1) {\textbf{I.}~Orientation};
  \node[pt, below left=of p1, xshift=-6mm] (p2) {\textbf{II.}~Field Theory};
  \node[pt, below right=of p1, xshift=6mm] (p3) {\textbf{III.}~The Calculus};
  \node[pt, below=18mm of p1] (p4) {\textbf{IV.}~Numerical Machinery};
  \node[pt, below=of p4] (p5) {\textbf{V.}~The Modalities};
  \node[pt, below=of p5] (p6) {\textbf{VI.}~The Collection};
  \node[ap, right=22mm of p4] (app) {Appendices:\\\footnotesize deep semantics};
  \draw[depedge] (p1) -- (p2);
  \draw[depedge] (p1) -- (p3);
  \draw[depedge] (p2) -- (p4);
  \draw[depedge] (p3) -- (p4);
  \draw[depedge] (p4) -- (p5);
  \draw[depedge] (p5) -- (p6);
  \draw[refedge] (app) -- (p3);
  \draw[refedge] (app) -- (p4);
  \node[cornertag, right=1mm of p1] {start here};
\end{tikzpicture}
\caption[Reading roadmap]{The parts and their dependencies. Solid arrows mean
``rests on''; the dotted arrows mark the appendices, which deepen the calculus
(Part~III) and the machinery (Part~IV) and can be read on demand.}
\label{fig:roadmap}
\end{figure}

\textbf{Part~I (Orientation)} is the map of the whole territory. It introduces
the six analyses, the program that runs them, and the single idea---fields as
tensors, steps as pure functions---that the rest of the book elaborates. Read it
first and in full; everything later refers back to it.

\textbf{Parts~II and~III} are the two halves of the foundation, and they are
independent of each other. \textbf{Part~II (The Field Theory)} builds the
physics and the finite-element method from your vector calculus.
\textbf{Part~III (The Abstract Framework)} builds the small functional calculus
from your linear algebra. A reader who wants the physics first can read
II~then~III; a reader who wants the computational language first can read
III~then~II. Both are needed before Part~IV.

\textbf{Part~IV (The Numerical Machinery)} is the algorithmic heart: how linear
systems, eigenproblems, time evolution, mesh adaptation, and physical-quantity
extraction actually run. \textbf{Part~V (The Constituent Modalities)} then walks
through the six analyses one at a time, each as a specific composition of the
machinery. \textbf{Part~VI (The Collection)} reassembles everything into the
synthesized library and steps back to view the whole machine.

The \textbf{appendices} hold the precise version of the computational
semantics---the formal grammar, the reduction and typing rules, the algebraic
laws, the floating-point conventions, and the inner details of the matrix-free
kernel. They are written with the same diagrams and discussion as the main text,
but they are reference material: reach for them when you want the exact rule
behind a picture, not on a first pass.

\section*{The order in which ideas are introduced}

Part~III teaches the functional calculus in a deliberate sequence, each idea
built only from the ones before it. If you have never programmed in this style,
follow this order and resist the urge to skip:

\begin{enumerate}[nosep]
  \item \emph{Values and immutability}---a tensor is a value you never modify.
  \item \emph{Higher-order functions}---functions that take and return functions.
  \item \emph{The fold}---one loop combinator, \texttt{iterate\_while}, from
  which every iteration is built.
  \item \emph{Records}---bundles of named fields, updated by making fresh copies.
  \item \emph{The state monad}---a disciplined way to thread one evolving piece
  of state so that its single point of change stays visible.
  \item \emph{Operators as values}---built once, applied many times.
  \item \emph{Shapes}---symbolic contracts on a tensor's axes.
  \item \emph{Rotations}---the idea that ties the layered representations together.
\end{enumerate}

\section*{The boxes}

Four kinds of boxed aside recur throughout, each with its own color so you can
find them at a glance:

\begin{keyidea}
A \textbf{Key idea} box distills the single most important takeaway of a section
into one or two sentences. If you remember nothing else from a section, remember
the box.
\end{keyidea}

\begin{pitfall}
A \textbf{Pitfall} box flags a tempting mistake---an algorithm that looks
correct but quietly computes the wrong thing, or a notation that means something
other than it appears to.
\end{pitfall}

\begin{physicsbox}
A \textbf{physical picture} box anchors an abstract construction to the concrete
device or experiment it models, so the mathematics never floats free of the
engineering.
\end{physicsbox}

\begin{notationbox}
A \textbf{Notation} box introduces or clarifies a piece of the pseudo-language
or the mathematical notation at the moment you first need it.
\end{notationbox}

\noindent Numbered \textbf{worked examples} carry a calculation all the way
through, by hand, on a small enough problem that you can check every step. They
are the rungs of the ladder; do them with pencil in hand.
